\documentclass[12pt,openright,openany,oneside,article,a4paper,brazi]{abntex2}
\usepackage[T1]{fontenc}
\usepackage[utf8]{inputenc}
\usepackage[brazil]{babel}           % estas quatro linhas
\usepackage[alf]{abntex2cite}
%\usepackage[brazil,brazilian]{babel}
\usepackage{lmodern,latexsym,enumerate}
\usepackage{amsmath,amsthm,amsopn,amstext,amscd,amsfonts,amssymb}
\usepackage[dvips]{graphicx}
\usepackage{setspace}
\usepackage{colortbl}
%\usepackage{natbib}
%\usepackage[loose]{subfigure}
%\usepackage{psfrag}
\usepackage{float}
\usepackage{array} % for defining a new column type
% \usepackage{varwidth}
 \usepackage{bm}
\usepackage{url}
\usepackage{indentfirst}
\usepackage{bbding}
\usepackage{pifont}
\usepackage{wasysym}
\usepackage{here,color,xcolor}
%\usepackage{abnt-alf}
%\usepackage[alf]{abntcite}
\usepackage{hyperref}
\usepackage[most]{tcolorbox}


\newcommand\signature[2]{% Nome; funcao
\begin{center}\begin{minipage}{10cm}
    \centering
    \vspace{14pt}\par
    \rule{10cm}{1pt}\par
   \large \textbf{#1}\par
    #2%
\end{minipage}
\end{center}
}


\newcommand\insertdate[1][\today]{
\vfill
\begin{center}
\large #1, \today
\end{center}
}


%\renewcommand{\familydefault}{phv}
\setlrmarginsandblock{2cm}{2cm}{*}
\setulmarginsandblock{2cm}{2cm}{*}
\checkandfixthelayout

\setlength{\ABNTEXcitacaorecuo}{1.8cm}
%configurações do formato do projeto
%\singlespacing Para um espaçamento simples
%\onehalfspacing %Para um espaçamento de 1,5
%    \doublespacing Para um espaçamento duplo
\usepackage{times} % times new roman
% pacote babel ou equivalent
%\usepackage[tmargin=2cm,bmargin=2cm,lmargin=2cm,rmargin=2cm]{geometry} % controle das margens 
\setlength{\parindent}{1.3cm}
\setlength{\parskip}{0.2cm} 
\usepackage{titlesec}

\titleformat{\section}[display]
{\large\bfseries}
{\thesection.}{}{}

\newtcolorbox{mybox1}[1]{colback=red!5!white,colframe=red!75!black,fonttitle=\bfseries,title=#1}

\newtcolorbox{mybox}{colback=gray!25,colframe=gray!25, sharp corners=all, height=0.80cm, valign=center}

% Definindo novas cores
\definecolor{verde}{rgb}{0,0.5,0}
% Configurando layout para mostrar codigos C++, R,...
\usepackage{listings}
\lstset{
  language=C++,
  basicstyle=\ttfamily\small,
  keywordstyle=\color{blue},
  stringstyle=\color{verde},
  commentstyle=\color{red},
  extendedchars=true,
  showspaces=false,
  showstringspaces=false,
  numbers=left,
  numberstyle=\tiny,
  breaklines=true,
  backgroundcolor=\color{green!10},
  breakautoindent=true,
  captionpos=b,
  xleftmargin=0pt,
}

\renewcommand{\bibsection}{\section*{Referências}}


\begin{document}

\begin{center}

\includegraphics[width=2.17cm,height=2.14cm]{ufcg.png}


{\textbf {UNIVERSIDADE FEDERAL DE CAMPINA GRANDE\\
PRÓ-REITORIA DE PESQUISA E EXTENSÃO\\
COORDENAÇÃO DE PESQUISA}}
\end{center}

 
\vspace{24pt}
\begin{mybox}
\begin{center}
{\Large \textbf{RELATÓRIO DE ATIVIDADES DO ALUNO}}
\end{center}
\end{mybox}
\vspace{12pt}
%%%%%%%%%%%%%%%%%%%%%%%%%%%%%%%%%%%%%%%%%%
%Informações do Projeto
%%%%%%%%%%%%%%%%%%%%%%%%%%%%%%%%%%%%%%%%%%


\begin{flushleft}
{\large \textbf{Programa:} PIBIC}\\
\vspace{14pt}
{\large \textbf{Título do Projeto:} Meu Projeto}\\
\vspace{14pt}
{\large \textbf{Aluno(a):} Nome do(a) Aluno(a)}\\
\vspace{14pt}
{\large \textbf{Orientador(a):} Nome do(a) orientador(a) }
\end{flushleft}
\vspace{200pt}

%%%%%%%%%%%%%%%%%%%%%%%%%%%%%%%%%%%%%%
%Informações Assinaturas
%%%%%%%%%%%%%%%%%%%%%%%%%%%%%%%%%%%%%%
\signature{Sr(a). Bolsista}{Bolsista}
\signature{Dr(a). Orientador(a)}{Orientador(a)}

\insertdate[Sumé]


\newpage

\pagenumbering{gobble}

\section*{Introdução}

 Em~\citeonline{RCT2017}
existe uma coleção de manuais completos sobre questões relacionadas a instalação e utilização do \textbf{R}, ver em \url{https://cran.r-project.org/manuals.html}. É possível encontrar vários livros com aplicações da linguagem \textbf{R} em diferente área, por exemplo, ver~\citeonline{Paradis2004,Baddeley2004, Morandat2012}.
Pacotes \textbf{R} são compostos de bibliotecas com funções e dados.


\section*{Objetivos}

\section*{Material e Métodos/Metodologia}

\section*{Resultados Parciais}

\newpage
\bibliographystyle{abntex2-alf}
\bibliography{ref}


\end{document}